
\exercise
Suppose $V$ is an inner product space and $v, w \in V\1$.  Define an operator
$T \in \L(V)$ by $Tu = \ip{u, v} w$.  Find a formula for~$\tr T\3$.


\exercise
Suppose $P \in \L(V)$ satisfies $P^2 = P\1$. Prove that
\[
\tr P = \dim \ran P\1.
\]
\exercisedisplayend


\exercise
Suppose $T \in \L(V)$ and $T^5 = T\3$. Prove that the real and imaginary parts of $\tr T$ are both integers.


\exercise
Suppose $V$ is an inner product space and $T \in \L(V)$. Prove that \label{393}
\[
\tr T^* = \overline{\tr T}.
\]
\exercisedisplayend


\exercise
Suppose $V$ is an inner product space.  Suppose $T \in \L(V)$ is a positive
operator and $\tr T = 0$. Prove that $T = 0$. \label{391}


\exercise
Suppose $V$ is an inner product space and $P\1, Q \in \L(V)$ are orthogonal projections. Prove that $\tr(PQ) \ge 0$.


\exercise
Suppose $T \in \L\paren[\big]{\C{3}}$ is the operator whose matrix is
\[
\mat{ccc}{
51 &-12 &-21 \\
60 &-40 &-28 \\
57 &-68 &1}.
\]
Someone tells you (accurately) that $-48$ and $24$ are eigenvalues of~$T\3$.
Without using a computer or writing anything down, find the third eigenvalue
of~$T\3$.


\exercise
Prove or give a counterexample: If $S, T \in \L(V)$, then
$
\tr (ST) = (\tr S) (\tr T)$.
%\exercisedisplayend


\exercise
Suppose $T \in \L(V)$ is such that $\tr(ST) = 0$ for all $S \in \L(V)$. Prove that $T = 0$.


\exercise
Prove that the trace is the only linear functional $\fun \tau {\L(V)} {\F{}}$ such that\label{9trChar}
\[
\tau(ST) = \tau(TS)\]
for all $S, T \in \L(V)$ and $\tau(I) = \dim V\1$.
\hint{Suppose that\/ $v_1, \ldots, v_n$ is a basis of\/ $V\1$. For\/ $j, k \in \br{1, \ldots, n}$, define\/ $P_{j,k} \in \L(V)$ by
$
P_{j,k}(a_1 v_1 + \cdots + a_n v_n) = a_k v_j$.
Prove that
\vspace{-7bp}
\[
\tau(P_{j,k}) = \begin{cases}
1 & \text{if } j = k, \\
0 &\text{if } j \ne k.
\end{cases}
\vspace{-6bp}
\]
Then for\/ $T \in \L(V)$, use the equation\/ $T = \sum_{k=1}^n \sum_{\js}^n \M(T)_{j,k} P_{j,k}$ to show that\/ $\tau(T) = \tr T$.}


\exercise
Suppose $V$ and $W$ are inner product spaces and $T \in \L(V\1,W)$.
Prove that if $e_1, \dots, e_n$ is an orthonormal basis of $V$ and $f_1, \ldots, f_m$ is an orthonormal basis of $W\3$, then\label{390}
\[
\tr\paren[\big]{T^*T} = \sum_{k=1}^n \sum_{\js}^m \abs{ \ip{Te_k, f_j} }^2\1.
\]
\exercisecomment{The numbers\/ $\ip{Te_k, f_j}$ are the entries of the matrix of\/ $T$ with respect to the orthonormal bases\/
$e_1, \dots, e_n$ and\/ $f_1, \ldots, f_m$. These numbers depend on the bases, but\/ $\tr (T^*T)$ does not depend on a choice of bases. Thus this exercise shows that the sum of the squares of the absolute values of the matrix entries does not depend on which orthonormal bases are used.}


\exercise
Suppose $V$ and $W$ are finite-dimensional inner product spaces.
\begin{subtheorem}
\item
Prove that
$
\ip{S, T} = \tr \paren[\big]{T^*S}
$
defines an inner product on $\L(V\1, W)$.
\item
Suppose $e_1, \dots, e_n$ is an orthonormal basis of $V$ and $f_1, \ldots, f_m$ is an orthonormal basis of $W\3$. Show that the inner product on $\L(V\1, W)$ from (a) is the same as the standard inner product on $\FF{mn}$, where we
identify each element of $\L(V\1, W)$ with its matrix (with respect to the bases just mentioned) and then with an element of
$\FF {mn}$.\index{Frobenius norm}\index{Hilbert--Schmidt norm}
\end{subtheorem}
\exercisecomment{Caution: The norm of a linear map\/ $T \in \L(V\1, W)$ as defined by \ref{7norm} is not the same as the norm that comes from the inner product in \uppar{a} above. Unless explicitly stated otherwise, always assume that\/ $\norm{T}$ refers to the norm as defined by \ref{7norm}.  The norm that comes from the inner product in \uppar{a} is called the \emph{Frobenius norm} or the \emph{Hilbert--Schmidt norm}.}


\exercise
Find $S, T \in \L\paren[\big]{ \P(\F{}) }$ such that $ST - TS = I$.\label{9commutator}
\hint{Make an appropriate modification of the operators in Example \ref{3noncom}.}
\exercisecomment{This exercise shows that additional hypotheses are needed on $S$ and $T$ to extend \ref{10ST} to the setting of infinite-dimensional vector spaces.}


